\chapter{Results}

\section*{Summary}
This chapter details the results of what was learned and accomplished during the project.

\section{Development Anticipated}
\label{sec:results-anticipated}

Based on the architecture in Chapter~4 and the methodology in Chapter~5, the following development outcomes are planned:

\textbf{Go Networking Daemon.} A working daemon that discovers peers on the local subnet through mDNS/DNS-SD, maintains persistent TCP connections with them, exchanges file metadata to detect changes, coordinates file transfers, and communicates with the C++ daemon through Unix domain sockets. The core packages (\texttt{discovery}, \texttt{network}, \texttt{ipc}) are already prototyped and tested independently; integration is the remaining work.

\textbf{C++ Storage Daemon.} A working daemon that watches the synchronization directory using kernel-level APIs (inotify on Linux, FSEvents on macOS), computes SHA-256 hashes of modified files with memory-mapped I/O, saves existing files to a \texttt{.versions/} directory before any overwrite, and streams incoming file data directly from TCP sockets to disk.

\textbf{Java Frontend.} A desktop application that shows the current synchronization status (which peers are connected, which files are syncing), displays a per-file version history with the ability to restore a previous version, and lets users configure the synchronization directory.

\textbf{Hypothesis Testing.} Two sets of measurements:
\begin{itemize}
    \item $H_1$ (Latency): timed synchronization tests across 3 to 5 peers on the campus WLAN, for files ranging from 1\,KB to 10\,MB. Target: sub-second average latency.
    \item $H_2$ (Consistency): concurrent edit scenarios where two peers modify the same file while disconnected, then reconnect. Expected outcome: automatic resolution with no permanent data loss, convergence within 30 seconds, and full recoverability from backups.
\end{itemize}

\section{Development Updated}
\label{sec:results-updated}

\textit{This section will be completed at the end of Action Learning Week in a separate standalone document. It will describe the differences between what was anticipated above and what was actually developed: features that were completed, features that were modified, and features that were deferred. Maximum one-half page.}
